%\ifodd\value{page}\hbox{}\newpage\fi
\chapter*{Śrī Lalitā Sahasranāma — Śloka 7}
\addcontentsline{toc}{chapter}{Śloka 7 - Nomi 16,17}

\begin{sloka}
{\sanskritfont
\begin{tabular}{c}
नवचम्पकपुष्पाभनासादण्डविराजिता ।\\
ताराकान्तितिरस्कारिनासाभरणभासुरा ॥ ७ ॥
\end{tabular}

\vspace{1.5em}

{\middlesize \iastfont
nava-campaka-puṣpābha-nāsā-daṇḍa-virājitā |\\
tārā-kānti-tiraskāri-nāsābharaṇa-bhāsurā ||7||}

\vspace{1em}

{\middlesize
\anvaya{%
नवचम्पकपुष्पाभनासादण्डविराजिता ।
ताराकान्तितिरस्कारिनासाभरणभासुरा ।
}
}

\middlesize \iastfont \itmean{%
«Il cui dorso del naso risplende come un fiore di \textit{campaka} appena sbocciato;  
il cui ornamento nasale brilla superando lo splendore delle stelle.»%
}

}
\end{sloka}

\wordmeaning{%
\wordline{नव-चम्पक-पुष्प-आभ-नासा-दण्ड-विराजिता}{nava-campaka-puṣpābha-nāsā-\\
daṇḍa-virājitā}
  { il cui dorso del naso risplende con il colore di un fiore di \textit{campaka} appena sbocciato;}%
\wordline{तारा-कान्ति-तिरस्कारि-नासा-आभरण-भासुरा}{tārā-kānti-tiraskāri-\\
nāsābharaṇa-bhāsurā}
  { il cui ornamento del naso è luminoso al punto da superare la luce delle stelle.}%
}

Lo \textit{Śloka} 7 prosegue la discesa contemplativa lungo il volto divino, soffermandosi sul naso come punto di equilibrio sottile e di discernimento. Il \textit{nāsā-daṇḍa}, paragonato a un fiore di \textit{campaka} appena sbocciato, non evoca soltanto delicatezza formale, ma freschezza, proporzione organica e armonia naturale. Il naso appare come asse verticale del volto, mediatore silenzioso tra lo sguardo e la parola.

La seconda metà dello \textit{śloka} introduce l’ornamento nasale (\textit{nāsābharaṇa}) non come semplice decorazione, ma come punto di condensazione luminosa. Il suo splendore, detto capace di superare la luce delle stelle, non rimanda a una luminosità cosmica dispersa, ma a una luce raccolta, prossima alla presenza cosciente stessa.

Dal punto di vista simbolico, il naso è il luogo in cui il respiro incontra la percezione. È il punto in cui il flusso vitale si fa ritmo consapevole, dove inspirazione ed espirazione si dispongono spontaneamente in equilibrio. Collocare qui una tale intensità luminosa significa indicare che la chiarezza nasce dalla misura, non dall’eccesso.

In questo \textit{śloka}, la bellezza di \textit{Lalitā} si manifesta come equilibrio assiale e sobrietà luminosa. Non è ornamento ridondante, ma forma centrata; non dispersione, ma presenza. La sua armonia rivela che la vera luce emerge là dove forma, respiro e coscienza coincidono.

Questo \textit{śloka} contiene due \textit{nāma} propriamente detti, che descrivono la forma divina secondo un principio di equilibrio assiale. \textbf{\textit{Nava-campaka-puṣpābha-nāsā-daṇḍa-virājitā}} qualifica il dorso del naso come segno di freschezza e proporzione naturale, indicando un’armonia che non deriva dall’ornamento, ma dalla misura intrinseca della forma. \textbf{\textit{Tārā-kānti-tiraskāri-nāsābharaṇa-bhāsurā}} specifica invece la luminosità dell’ornamento nasale, non come brillantezza esteriore, ma come luce concentrata che nasce dalla prossimità alla coscienza.

